\chapter{Как устроена контейнеризация}

Контейнер кажется магией только до тех пор, пока не увидишь, что под ним.
Никакой эмуляции компьютера нет: ядро Linux умеет изолировать группы
процессов, и Docker всего лишь удобно этим пользуется. В этой главе
разобраны четыре механизма~--- пространства имён, контрольные группы,
слоистая файловая система и клиент-серверная архитектура Docker,~---
а затем прослежен путь команды \cmd{docker run} от терминала до процесса.

\section{Пространства имён (namespaces)}

Пространство имён~--- механизм ядра, который даёт процессу собственный
взгляд на некоторый глобальный ресурс. Процессы в разных пространствах
имён видят разные файловые системы, разные сетевые интерфейсы, разные
списки процессов, хотя работают на одном ядре.

\begin{table}[htbp]
\centering
\small
\begin{tabularx}{\textwidth}{@{}lXX@{}}
\toprule
\textbf{Namespace} & \textbf{Что изолирует} & \textbf{Как это заметно} \\
\midrule
\cmd{pid} & номера процессов & внутри контейнера \cmd{ps} видит только
процессы контейнера, главный из них имеет PID~1 \\
\cmd{net} & интерфейсы, порты, маршруты, \cmd{iptables} & свой \cmd{localhost},
свои порты, свой DNS \\
\cmd{mnt} & точки монтирования & своя корневая файловая система \\
\cmd{uts} & имя хоста и домен & \cmd{hostname} внутри не совпадает с хостом \\
\cmd{ipc} & разделяемая память, семафоры, очереди & свои средства
межпроцессного взаимодействия \\
\cmd{user} & идентификаторы UID/GID & \cmd{root} внутри может быть обычным
пользователем снаружи \\
\cmd{cgroup} & видимая иерархия контрольных групп & контейнер не видит
чужие лимиты \\
\cmd{time} & системные часы & можно задать смещение времени (новый
namespace в свежих ядрах) \\
\bottomrule
\end{tabularx}
\caption{Пространства имён Linux, которые использует контейнер}\label{tab:namespaces}
\end{table}

Docker создаёт для контейнера набор пространств имён (кроме \cmd{user},
если не включён rootless-режим или \cmd{userns-remap}). С точки зрения
приложения оно выглядит как отдельная машина: свой хостнейм, свой
\cmd{localhost}, свой корень файловой системы.

\begin{notebox}
Изоляцию можно увидеть без Docker. Команда
\cmd{unshare --pid --fork --mount-proc ps aux} запускает \cmd{ps}
в новом пространстве имён процессов: вы увидите только саму команду,
а не сотни процессов хоста. Для этой команды нужны права root
или включённые user namespaces.
\end{notebox}

\begin{warnbox}
PID~1 внутри контейнера~--- это ваше приложение, а не \cmd{systemd}
и не \cmd{init}. Оно обязано обрабатывать сигналы и \enquote{пожинать}
зомби-процессы. Если приложение не умеет быть лидером, используйте
флаг \flag{--init}: Docker добавит маленький init-процесс, который
корректно передаёт сигналы и убирает зомби.
\end{warnbox}

\section{Контрольные группы (cgroups)}

Пространства имён отвечают за \emph{что видно}, а контрольные группы~---
за \emph{сколько можно}. Cgroups ограничивают и учитывают потребление
ресурсов: процессорного времени, памяти, дискового ввода-вывода,
числа процессов.

Docker выставляет эти ограничения флагами команды \cmd{docker run}
или в compose-файле:

\begin{code}[language=bash]
$ docker run --memory=512m --cpus=1.5 --pids-limit=100 nginx:alpine
\end{code}

\begin{center}
\small
\begin{tabularx}{\textwidth}{@{}lXl@{}}
\toprule
\textbf{Флаг} & \textbf{Ограничение} & \textbf{Пример} \\
\midrule
\flag{--memory} & максимум оперативной памяти; при превышении процесс
убивается (код 137) & \flag{--memory=512m} \\
\flag{--memory-swap} & суммарный лимит памяти и swap & \flag{--memory-swap=1g} \\
\flag{--cpus} & доля процессорного времени & \flag{--cpus=1.5} \\
\flag{--cpu-shares} & относительный приоритет при конкуренции & \flag{--cpu-shares=512} \\
\flag{--pids-limit} & максимальное число процессов и потоков & \flag{--pids-limit=200} \\
\flag{--blkio-weight} & относительный приоритет дискового ввода-вывода & \flag{--blkio-weight=300} \\
\bottomrule
\end{tabularx}
\end{center}

Современные системы используют cgroups~v2 с единой иерархией; старые~---
cgroups~v1 с отдельными контроллерами. Для пользователя Docker разница
почти незаметна, но при отладке лимитов полезно знать, что смотреть нужно
в \file{/sys/fs/cgroup}.

\begin{tipbox}
Лимит памяти~--- не украшение. Без \flag{--memory} один контейнер
с утечкой памяти способен исчерпать память хоста и утащить за собой
все остальные процессы. В production лимиты задавайте всегда, даже
приблизительные.
\end{tipbox}

\section{Слоистая файловая система}

Образ не копируется целиком при каждом запуске. Он хранится набором
слоёв, доступных только для чтения, а драйвер \cmd{overlay2} собирает
из них единое дерево. Сверху контейнер получает один записываемый слой.

\begin{itemize}
  \item \textbf{Чтение} идёт из нижних слоёв без изменений.
  \item \textbf{Запись} в существующий файл вызывает copy-up: файл
        копируется в верхний слой, и изменение применяется к копии.
  \item \textbf{Удаление} файла из нижнего слоя создаёт белый маркер
        (whiteout): файл скрыт, но остаётся в образе.
  \item \textbf{Удаление контейнера} уничтожает его записываемый слой
        вместе со всеми изменениями.
\end{itemize}

Отсюда два важных следствия. Во-первых, образ можно безопасно использовать
для множества контейнеров: они не портят его слои. Во-вторых, данные,
которые должны пережить контейнер, нужно выносить в тома или привязки
(глава~7), а не хранить в записываемом слое.

\begin{notebox}
Слои разделяются между образами. Если десять образов основаны
на \cmd{ubuntu:24.04}, базовый слой хранится на диске один раз.
Поэтому \cmd{docker pull} второго похожего образа часто идёт быстро:
большая часть слоёв уже есть локально.
\end{notebox}

\section{Клиент, демон и среда выполнения}

Docker построен как клиент-серверная система. Команда \cmd{docker}
не запускает контейнеры сама: она отправляет запрос по HTTP демону
\cmd{dockerd}, который работает в фоне и делает всю работу.

\begin{figure}[htbp]
\centering
\begin{tikzpicture}[
  node distance=6mm and 7mm,
  font=\sffamily\scriptsize,
  every node/.style={align=center},
  component/.style={draw=accent!55!black, fill=accent!7, rounded corners=2pt,
    minimum width=2.35cm, minimum height=1.05cm, inner sep=3pt},
  proc/.style={draw=warm!70!black, fill=warm!8, rounded corners=2pt,
    minimum width=2.35cm, minimum height=1.05cm, inner sep=3pt},
  registry/.style={draw=accent2!70!black, fill=accent2!7, dashed,
    rounded corners=2pt, minimum width=2.6cm, minimum height=0.95cm, inner sep=3pt},
  arr/.style={-Latex, thick, accent!70},
  darr/.style={-Latex, thick, dashed, accent2!80!black},
]
  \node[registry] (reg) {реестр образов\\Docker Hub и др.};
  \node[component, below=14mm of reg] (cli) {\textbf{docker CLI}\\[1pt]\cmd{docker run}};
  \node[component, right=of cli] (dockerd) {\textbf{dockerd}\\[1pt]API, образы,\\сети, тома};
  \node[component, right=of dockerd] (containerd) {\textbf{containerd}\\[1pt]управление\\жизненным циклом};
  \node[component, right=of containerd] (shim) {\textbf{containerd-shim}\\[1pt]держит stdio,\\собирает код выхода};
  \node[proc, right=of shim] (proc) {\textbf{процесс}\\[1pt]PID~1 в контейнере};

  \draw[arr] (cli) -- node[above=1pt, text=muted] {HTTP} node[below=1pt, text=muted] {Unix-сокет} (dockerd);
  \draw[arr] (dockerd) -- node[above=1pt, text=muted] {gRPC} (containerd);
  \draw[arr] (containerd) -- node[above=1pt, text=muted] {задача} (shim);
  \draw[arr] (shim) -- node[above=1pt, text=muted] {runc} (proc);
  \draw[darr] (reg) -- node[right=2pt, text=muted] {pull/push} (dockerd);
\end{tikzpicture}
\caption{Путь команды \cmd{docker run}: клиент обращается к демону, тот
передаёт работу containerd, а процесс создаётся низкоуровневой средой
выполнения}\label{fig:architecture}
\end{figure}

Разберём элементы схемы.

\begin{description}[style=nextline]
  \item[docker CLI] Клиент командной строки. Ничего не знает о внутреннем
        устройстве контейнеров: просто формирует HTTP-запросы и печатает
        ответы. Может работать с удалённым демоном.
  \item[dockerd] Демон Docker. Принимает API-запросы, управляет образами,
        томами, сетями, а создание контейнеров делегирует containerd.
  \item[containerd] Менеджер контейнеров: скачивает образы, распаковывает
        слои, запускает задачи, следит за жизненным циклом. Именно его
        использует Kubernetes, поэтому контейнеры Docker и Kubernetes
        совместимы на уровне образов.
  \item[containerd-shim] Небольшой процесс-посредник. Позволяет демону
        перезапускаться, не убивая контейнеры, удерживает потоки
        ввода-вывода и сообщает код выхода.
  \item[runc] Реализация стандарта OCI Runtime. Получает bundle с корневой
        файловой системой и конфигурацией, вызывает \cmd{clone()}
        с нужными namespaces, настраивает cgroups, монтирует rootfs
        и выполняет \cmd{exec} команды. После запуска процесса runc
        завершается~--- контейнер живёт без него.
\end{description}

\begin{notebox}
Docker Desktop для macOS и Windows запускает \cmd{dockerd} внутри
лёгкой виртуальной машины, потому что ядро этих систем не является
Linux. С точки зрения CLI всё выглядит так же, но путь к сокету
и производительность файловой системы отличаются от Linux.
\end{notebox}

\section{Что происходит при \cmd{docker run nginx:alpine}}

Полезно один раз проследить всю цепочку целиком.

\begin{enumerate}
  \item CLI разбирает аргументы и отправляет демону запрос
        \cmd{POST /containers/create}.
  \item Демон проверяет, есть ли образ \cmd{nginx:alpine} локально.
        Если нет~--- обращается к реестру, скачивает манифест и слои,
        проверяя контрольные суммы.
  \item Слои распаковываются в хранилище образов (snapshotter);
        из них собирается корневая файловая система контейнера.
  \item Создаются метаданные контейнера, выделяется сетевой адрес,
        подготавливаются точки монтирования томов и конфигурация.
  \item containerd получает задачу, запускает shim, а тот вызывает runc.
  \item runc создаёт пространства имён и cgroups, делает \cmd{pivot\_root}
        в файловую систему образа и запускает процесс \cmd{nginx}.
  \item runc завершается; shim остаётся, удерживая stdio и ожидая
        код выхода. Процесс \cmd{nginx} становится PID~1 внутри контейнера.
\end{enumerate}

\begin{notebox}
Команда \cmd{docker run}~--- это сокращение для \cmd{docker create}
и \cmd{docker start}. Если нужно подготовить контейнер, но не запускать
его, используйте \cmd{docker create}; подробнее~--- в главе~5.
\end{notebox}

\section{Куда CLI отправляет команды}

По умолчанию клиент подключается к локальному демону.
Канал связи~--- Unix-сокет \file{/var/run/docker.sock}.
Это поведение можно изменить:

\begin{code}[language=bash]
$ docker context ls
$ docker context create remote --docker "host=ssh://user@server"
$ docker context use remote
$ docker --context default ps
$ export DOCKER_HOST=ssh://user@server
\end{code}

Контексты удобны, когда нужно переключаться между локальным Docker,
удалённым сервером и кластером. Переменная \envvar{DOCKER\_HOST}~
более старый способ; контекст имеет приоритет и хранит дополнительные
настройки.

\begin{warnbox}
Доступ к \file{/var/run/docker.sock} равносилен доступу root на хосте:
через демон можно смонтировать корень хоста внутрь контейнера.
Не добавляйте посторонних пользователей в группу \cmd{docker}
на общих машинах и не пробрасывайте сокет в недоверенные контейнеры.
Подробнее~--- в главе~12.
\end{warnbox}

\begin{summarybox}
\begin{itemize}
  \item Контейнер~--- это процесс с набором namespaces и лимитами cgroups.
  \item Namespaces изолируют видимость ресурсов, cgroups~--- потребление.
  \item Образы состоят из неизменяемых слоёв; контейнер получает
        один записываемый слой, который исчезает при удалении.
  \item Docker~--- клиент-серверная система: CLI общается с \cmd{dockerd},
        тот делегирует containerd, а процесс создаёт runc.
  \item Сокет Docker~--- мощный и опасный интерфейс; доступ к нему нужно
        контролировать.
\end{itemize}
\end{summarybox}

\section*{Вопросы для самопроверки}
\begin{selfcheck}
  \item Какие пространства имён создаёт Docker и что изолирует каждое?
  \item Чем namespace отличается от cgroup?
  \item Что такое copy-up и whiteout в overlayfs?
  \item Почему runc завершается, а контейнер продолжает работать?
  \item Что делает containerd-shim и зачем он нужен?
  \item Чем опасен смонтированный в контейнер \file{/var/run/docker.sock}?
\end{selfcheck}
